\section{Verifikation und Validierung}
\label{sec:validation}

Dieses Kapitel beschreibt die Verifikation und Validierung der
prototypischen Realisierung. Untersucht werden die regelbasiert prüfbaren
Implementierungseigenschaften, das Zusammenspiel der Verarbeitungsschritte,
die durchgängige Nachverfolgbarkeit sowie die Validierung und menschliche
Freigabe des erzeugten SysML-v2-Artefakts. Die Ergebnisse dieser Prüfungen
werden in Kapitel~\ref{sec:results} dargestellt.

\subsection{Ziel und Aufbau der Evaluation}
\label{sec:validation_objectives}

Die durchgeführten Tests des Prototyps konzentrierten sich auf zwei Ziele.
Erstens wurde untersucht, ob sich das in dieser Arbeit entwickelte
Architekturkonzept für die Transformation dokumentenbasierter
Systembeschreibungen in formale SysML-v2-Repräsentationen innerhalb des
gewählten Untersuchungsrahmens technisch umsetzen und durchgängig anwenden
lässt. Zweitens wurde geprüft, ob die darin definierten Konzepte und
Transformationsschritte im Prototyp entsprechend ihrer Spezifikation
implementiert wurden.

Für die Untersuchung wurden ausgewählte Systembeschreibungen bestehender
Produkte der Firma PreciPoint herangezogen. Die eingeschränkte Funktionalität
des Prototyps wurde bei der Auswahl und Beurteilung der Testfälle
berücksichtigt
(siehe Tabelle~\ref{tab:implementation_limits}). Die systematische
Untersuchung der Eigenschaften, des Verhaltens und der Grenzen eines
entwickelten Artefakts entspricht dem gestaltungsorientierten
Forschungsansatz dieser Arbeit
\cite{Wieringa2014,Peffers2007,ISO42030}.

Die Evaluation verbindet mehrere einander ergänzende Prüfbereiche:

\begin{itemize}
    \item Die automatisierte Software- und Vertragsverifikation untersucht
          regelbasiert überprüfbare Eigenschaften der Implementierung. Dazu
          gehören unter anderem Datenstrukturen, Identitäten,
          Persistenzregeln, Zustandsübergänge und Provenienzbeziehungen.
          Zusätzlich wird geprüft, ob das in
          Abschnitt~\ref{sec:architecture_controlled_transformation}
          eingeführte \emph{fail-closed}-Prinzip technisch eingehalten wird.

    \item Die End-to-End-Verifikation untersucht den durchgängigen
          Verarbeitungspfad von den Ausgangsinformationen bis zum erzeugten
          SysML-v2-Artefakt. Dabei wird geprüft, ob die einzelnen
          Transformationsschritte im Zusammenspiel funktionieren, die
          vorgesehenen Autorisierungsgrenzen technisch eingehalten werden
          und die erforderlichen Herkunftsinformationen erhalten bleiben.

    \item Die Verarbeitung mehrerer Engineering-Quellen erweitert diese
          Untersuchung auf den projektweiten Kontext. Dabei wird geprüft,
          ob Quellen zunächst getrennt verarbeitet und anschließend
          gemeinsam berücksichtigt werden können, während ihre Identität,
          Provenienz und fachlichen Autorisierungszustände erhalten bleiben.

    \item Die technische Validierung des erzeugten SysML-v2-Artefakts
          untersucht dessen formale und technische Verarbeitbarkeit in der
          verwendeten SysML-v2-Toolumgebung. Ergänzend wird geprüft, ob der
          Prototyp die für menschliche Autorisierungsentscheidungen sowie die
          abschließende Modellprüfung und Freigabe erforderlichen
          Informations- und Modellzustände bereitstellt und die jeweiligen
          Entscheidungen eindeutig mit den geprüften Zuständen verknüpfen
          kann.
\end{itemize}

Die Verifikation des durchgängigen Verarbeitungspfads wurde formativ
durchgeführt. Während der Entwicklung wurde geprüft, ob die einzelnen
Verarbeitungsschritte und Autorisierungsgrenzen korrekt zusammenspielen.
Beobachtete Architektur- oder Implementierungsprobleme wurden konkreten
Verarbeitungsschritten zugeordnet, korrigiert und anschließend erneut
geprüft. Die formative Verifikation unterstützte damit zugleich die
schrittweise Konsolidierung des Prototyps. Für eindeutig spezifizierbare
Eigenschaften wurden soweit möglich automatisierte Modul-, Vertrags-,
Integrations- und Regressionstests eingesetzt. :contentReference[oaicite:2]{index=2}

Die einzelnen Prüfverfahren liefern unterschiedliche Aussagen über den
Prototyp. Bestandene Software- und Vertragstests bestätigen jeweils die
konkret spezifizierte und geprüfte technische Eigenschaft. Die
End-to-End- und Mehrquellenverifikation untersuchen das Zusammenspiel der
Verarbeitungsschritte und ihrer Übergänge. Die technische
SysML-v2-Validierung bestätigt die Verarbeitbarkeit des erzeugten
Zielartefakts in der verwendeten Toolumgebung. Die Ergebnisse werden in den
folgenden Abschnitten jeweils unmittelbar gemeinsam mit dem zugehörigen
Prüfverfahren und den verwendeten Akzeptanzkriterien dargestellt.

Nach Abschluss dieser Prüfungen wurde der konsolidierte Prototyp zusätzlich
in einer extern begleiteten Demonstration über seinen regulären
Verarbeitungspfad ausgeführt. Dabei entstand ein vollständig durchlaufener
und dauerhaft gespeicherter Referenzlauf. Dieser bildet die Grundlage für
die in Kapitel~\ref{sec:results} dargestellte exemplarische Rekonstruktion
der durchgängigen Nachverfolgbarkeit. :contentReference[oaicite:3]{index=3}

\subsection{Automatisierte Software- und Vertragsverifikation}
\label{sec:validation_software}

Die automatisierte Softwareverifikation untersucht die regelbasiert
überprüfbaren Eigenschaften des Prototyps. Im Mittelpunkt stehen Elemente,
die eindeutig spezifiziert sind und somit automatisiert geprüft werden
können. Dazu gehören beispielsweise Datenstrukturen, Identitäten,
Persistenzregeln, zulässige Zustandsübergänge und die Verknüpfung
aufeinanderfolgender Verarbeitungsschritte. Die gezielte Prüfung solcher
Elemente bildet einen Bestandteil der systematischen Evaluation des
entwickelten Prototyps
\cite{Wieringa2014,ISO42030}.

Zu den automatisiert überprüfbaren Elementen gehören auch die in
Kapitel~\ref{sec:implementation} eingeführten technischen Verträge zwischen
den Verarbeitungsschritten. Ein solcher Vertrag legt fest, welche
Eingangszustände ein Verarbeitungsschritt verwenden darf, welche Identitäten
und Referenzen erhalten bleiben müssen, welche Bedingungen für eine gültige
Verarbeitung gelten und welche Ergebniszustände weitergegeben werden können.
In der Implementierung werden diese Festlegungen als \emph{Contracts}
bezeichnet. Ihre automatisierte Prüfung kontrolliert, ob die einzelnen
Verarbeitungsschritte auf den dafür vorgesehenen Informations- und
Entscheidungszuständen aufbauen.

Die Prüfungen erfolgen auf drei Ebenen
(Tabelle~\ref{tab:validation_software_levels}). Auf der ersten Ebene werden
einzelne Funktionen, Datenstrukturen und technische Verträge geprüft. Die
zweite Ebene betrachtet das Zusammenspiel mehrerer Verarbeitungsschritte und
insbesondere die Übergänge zwischen ihnen. Die dritte Ebene umfasst die
wiederholte Ausführung des vollständigen automatisierten Testsatzes nach
Änderungen am Prototyp.

\begin{table}[htbp]
    \centering
    \small
    \caption[Ebenen der automatisierten Softwareverifikation]{
        Automatisierte Softwareverifikation von technischen Regeln,
        kontrollierten Übergängen und dem Gesamtsystem.}
    \label{tab:validation_software_levels}
    \begin{tabular}{p{3.6cm}p{5.0cm}p{4.8cm}}
        \hline
        \textbf{Prüfebene} &
        \textbf{Prüfgegenstand} &
        \textbf{Beispielhafte Kriterien} \\
        \hline

        Komponenten- und Vertragsprüfung &
        einzelne Funktionen, Datenstrukturen und definierte
        Übergabebedingungen &
        eindeutige Identitäten, Persistenz,
        Schema- und Typregeln sowie
        inhaltsgebundene Prüfwerte (\emph{Fingerprints}) \\

        Integrations- und Grenzprüfung &
        Zusammenspiel mehrerer Verarbeitungsschritte und Übergänge
        zwischen ihnen &
        erforderliche Autorisierungen,
        Erhalt der Provenienz,
        zulässige Zustandsübergänge und
        Einhaltung des \emph{fail-closed}-Prinzips \\

        Regression des Gesamtsystems &
        bereits implementierte und zuvor geprüfte Funktionen nach
        Änderungen am Prototyp &
        weiterhin erfolgreich ausführbare Verarbeitungs- und
        Autorisierungspfade \\
        \hline
    \end{tabular}
\end{table}

Neben den regulären Verarbeitungspfaden wurden gezielt Grenz- und
Fehlerfälle geprüft. Dazu gehörten beispielsweise fehlende erforderliche
Autorisierungen, widersprüchliche Identitäten, veränderte inhaltsgebundene
Prüfwerte, unvollständige Referenzen oder Modellierungszustände, für die
keine zulässige Weiterverarbeitung definiert ist. Dabei wurde geprüft, ob
die Verarbeitung entsprechend dem in
Abschnitt~\ref{sec:architecture_controlled_transformation} eingeführten
\emph{fail-closed}-Prinzip an der vorgesehenen Stelle blockiert wird.

Für KI-gestützte Verarbeitungsschritte konzentrierte sich die
automatisierte Verifikation auf technisch eindeutig spezifizierbare
Eigenschaften. Dazu gehörten insbesondere die Einhaltung vorgegebener
Datenstrukturen, die Bindung eines Verarbeitungsergebnisses an seine
Engineering-Quelle, vollständige Referenzen sowie die Trennung zwischen
maschinell erzeugten Interpretationsergebnissen und menschlich
autorisierter Engineering-Information. Damit wurde geprüft, ob die für die
spätere fachliche Bewertung erforderlichen Informationszustände und
Verknüpfungen entsprechend dem Architekturkonzept bereitgestellt werden.

Nach Änderungen wurden zunächst die unmittelbar betroffenen Funktionen und
Verträge geprüft. Anschließend wurden die angrenzenden
Verarbeitungsschritte gemeinsam getestet. Abschließend wurde der
vollständige automatisierte Testsatz des Prototyps erneut ausgeführt.
Ergänzende technische Konsistenzprüfungen kontrollierten den jeweiligen
Repository- und Implementierungsstand.

\paragraph{Ergebnisse}

Die automatisierte Software- und Vertragsverifikation bestätigte für die
jeweils geprüften Entwicklungsstände die spezifizierten technischen
Eigenschaften des Prototyps sowie die implementierten Verträge zwischen den
Verarbeitungsschritten. Geprüft wurden insbesondere Datenstrukturen,
Identitäten, Persistenzregeln, Zustandsübergänge sowie die dafür
festgelegten Referenz- und Autorisierungsbedingungen.

Für einen vor Abschluss der Verarbeitung mehrerer Quellen akzeptierten
Entwicklungsstand wurden 6100 automatisierte Tests erfolgreich ausgeführt.
Die nachfolgenden Änderungen im Zusammenhang mit der Verarbeitung mehrerer
Quellen wurden durch gezielte Prüfungen der betroffenen Funktionen und
Verträge sowie anschließende Regressionstests abgesichert. Die ergänzende
technische Konsistenzprüfung bestätigte dabei einen konsistenten
Repository-Zustand.

Die Anzahl erfolgreich ausgeführter Tests dient dabei als Nachweis für die
jeweils spezifizierten und geprüften technischen Eigenschaften. Die
Ergebnisse zeigen, dass diese Eigenschaften und Verträge in den geprüften
Entwicklungsständen erfolgreich ausgeführt und nach Änderungen erneut
verifiziert werden konnten. Das Zusammenspiel der einzelnen Mechanismen im
durchgängigen Transformationsprozess wird im folgenden Abschnitt anhand der
End-to-End-Verifikation untersucht.

\subsection{Verifikation des durchgängigen Verarbeitungspfads}
\label{sec:validation_e2e_multisource}

Ergänzend zur automatisierten Software- und Vertragsverifikation wurde der
Prototyp anhand ausgewählter realer Systembeschreibungen über
zusammenhängende Verarbeitungspfade geprüft. Ziel war festzustellen, ob die
einzelnen Verarbeitungsschritte, technischen Verträge und
Autorisierungsgrenzen im Zusammenspiel einen kontrollierten
Transformationsprozess bilden.

Für die formative Verifikation wurde zunächst ein vollständiger
Verarbeitungspfad von einer registrierten Engineering-Quelle bis zum
erzeugten SysML-v2-Artefakt über die reguläre Anwendungsschnittstelle des
Prototyps ausgeführt. Während der Verarbeitung wurden die für die
vorgesehenen menschlichen Entscheidungen erforderlichen
Informationszustände bereitgestellt und die entsprechenden
Autorisierungsentscheidungen durch den Autor eingegeben. An den Übergängen
zwischen den Verarbeitungsschritten wurde geprüft, ob die erforderlichen
Vorgängerzustände, Provenienzbeziehungen und Autorisierungsentscheidungen
vorlagen und korrekt miteinander verknüpft waren. Die weitere Verarbeitung
erfolgte jeweils auf Grundlage der für den Übergang festgelegten
Voraussetzungen.

Die Untersuchung war formativ angelegt. Beobachtete Abweichungen wurden dem
betroffenen Architekturmechanismus oder Verarbeitungsschritt zugeordnet und
dokumentiert. Nach einer Korrektur wurden zunächst die unmittelbar
betroffenen Softwarefunktionen und anschließend der entsprechende Abschnitt
des Verarbeitungspfads erneut geprüft. Frühere Fehlerzustände blieben dabei
als eigene Zustände erhalten und standen für die spätere Auswertung zur
Verfügung.

\paragraph{Ergebnisse der formativen Verifikation}

Die formative Verifikation des durchgängigen Verarbeitungspfads identifizierte
neben technischen Implementierungsfehlern mehrere Befunde mit unmittelbarer
Bedeutung für das Architekturkonzept. Als architekturrelevant wurden
Beobachtungen behandelt, bei denen eine vorgesehene Informations-,
Provenienz- oder Autorisierungsgrenze für die kontrollierte
Weiterverarbeitung präzisiert werden musste. Daraus wurden dokumentierte
Architekturentscheidungen sowie Anpassungen der technischen Verträge
abgeleitet. Die betroffenen Verarbeitungsschritte wurden anschließend erneut
geprüft.

Tabelle~\ref{tab:validation_architecture_findings} fasst die für die
Architekturentwicklung wesentlichen Befunde zusammen.

\begin{table}[htbp]
    \centering
    \small
    \caption[Architekturrelevante Befunde der formativen Verifikation]{
        Architekturrelevante Befunde der formativen Verifikation des
        Verarbeitungspfads und die daraus abgeleiteten Anpassungen.}
    \label{tab:validation_architecture_findings}
    \begin{tabular}{p{4.1cm}p{4.8cm}p{4.7cm}}
        \hline
        \textbf{Beobachtung} &
        \textbf{Architektonische Reaktion} &
        \textbf{Resultierender Zustand} \\
        \hline

        Verarbeitungskontext konnte als fachliche Information interpretiert
        werden &
        Trennung von Engineering-Quelle, Kontextquelle und
        Verarbeitungskontext &
        Engineering-Quellen bilden die fachliche Grundlage für
        quellengebundene Evidenz; Kontextinformationen und
        Verarbeitungskontext bleiben davon getrennt \\

        Persona-Verarbeitung erzeugte unterschiedliche Mengen fachlicher
        Gegenstände &
        Bildung gemeinsamer fachlicher Referenzgegenstände vor der
        Persona-Interpretation &
        alle Persona-Perspektiven interpretieren dieselben
        \emph{Engineering Subjects} \\

        akzeptierte fachliche Beziehungen fehlten im autorisierten
        Übergang zur Modellbildung &
        Einbeziehung autorisierter fachlicher Beziehungen in den
        weitergegebenen Informationszustand &
        autorisierte fachliche Gegenstände und Beziehungen stehen gemeinsam
        für die nachfolgende Modellbildung zur Verfügung \\

        fachliche Bedeutung wurde unmittelbar mit der Zielrepräsentation
        verknüpft &
        Trennung von fachlicher Bedeutung, struktureller Einordnung und
        zielmodellspezifischer Formulierung &
        fachliche Autorisierung und modellbezogene Autorisierung bilden
        getrennte Entscheidungsstufen \\

        während der nachgelagerten Modellbildung entstanden zusätzliche
        fachliche oder modellbezogene Entscheidungsbedarfe &
        Begrenzung der nachgelagerten Transformation auf eindeutig
        regelgebundene Übergänge &
        offene Entscheidungen werden an die dafür vorgesehene
        Autorisierungsgrenze zurückgeführt \\
        \hline
    \end{tabular}
\end{table}

Die Befunde führten damit während der prototypischen Realisierung zur
Präzisierung zentraler Architekturmechanismen. Dazu gehören die Trennung der
Informationsrollen, die gemeinsame Bildung fachlicher Referenzgegenstände,
die getrennte fachliche und modellbezogene Autorisierung sowie die
regelgebundene Transformation nach abgeschlossener Autorisierung. Die
resultierenden Mechanismen sind in
Kapitel~\ref{sec:architecture} beschrieben.

Technische oder bedienungsbezogene Befunde wurden ebenfalls korrigiert und
erneut geprüft. Die vollständige Historie der dokumentierten Befunde mit
Korrektur, Wiederholungsprüfung und abschließendem Status ist in
Anhang~\ref{app:findings_register} dokumentiert.

\paragraph{Verifikation der Verarbeitung mehrerer Quellen}

Zusätzlich wurde untersucht, ob mehrere Engineering-Quellen innerhalb
desselben Projekts zunächst getrennt verarbeitet und anschließend gemeinsam
für die Modellbildung berücksichtigt werden können. Hierfür wurden vier
aufbereitete und inhaltlich aufeinander bezogene Testdokumente verwendet.
Sie repräsentierten unterschiedliche Informationsperspektiven auf dasselbe
technische Projekt und wurden jeweils als eigenständige Engineering-Quelle
registriert.

Die Quellen wurden zunächst unabhängig voneinander verarbeitet und die
vorgesehenen fachlichen Autorisierungszustände erzeugt. Anschließend wurde
die projektweite Zusammenführung geprüft. Dabei wurden folgende
Akzeptanzkriterien verwendet:

\begin{itemize}
    \item Die projektbezogene Quellenzulassung (\emph{Project Fit}) wird bei
          der Entscheidung über die projektweite Verwendung einer
          Engineering-Quelle berücksichtigt.

    \item Jede Engineering-Quelle und der zugehörige Verarbeitungslauf
          bleiben eindeutig identifizierbar.

    \item Quellengebundene Evidenz, menschliche Prüfentscheidungen und
          autorisierte Engineering-Information bleiben mit ihrer jeweiligen
          Quelle und ihrem Verarbeitungslauf verknüpft.

    \item Die projektweite Verarbeitung verwendet fachlich autorisierte
          Engineering-Information aus entsprechend zugelassenen Quellen.

    \item Informationen mehrerer zugelassener Quellen können gemeinsam für
          die Modellbildung verwendet werden, während ihre jeweilige
          Herkunft erhalten bleibt.

    \item Lokal wiederverwendete technische Identifikatoren aus
          unterschiedlichen Verarbeitungsläufen bleiben durch ihren
          Quellen- und Verarbeitungskontext eindeutig unterscheidbar.

    \item Fachlich autorisierte Zustände entstehen über die dafür
          vorgesehenen menschlichen Autorisierungsentscheidungen und bleiben
          mit dem jeweils geprüften Informationszustand verknüpft.
\end{itemize}

Die Prüfung adressierte damit den Übergang von einzelnen quellenbezogenen
Verarbeitungspfaden zu einem gemeinsamen Projektkontext. Die projektweite
fachliche Bewertung wird durch die maschinelle Identifikation von
Ähnlichkeiten, Ergänzungen, möglichen Widersprüchen und Unsicherheiten
vorbereitet. Für die technische Verifikation wurde dabei insbesondere
geprüft, ob erforderliche Entscheidungsgrundlagen bereitgestellt, die
entstehenden Zustände nachvollziehbar verknüpft und die vorgesehenen
Verarbeitungsgrenzen eingehalten werden.

Bei einem Halt der Verarbeitung wurde anhand des gespeicherten Zustands
geprüft, welcher Auslöser vorlag. Ein durch eine fehlende Voraussetzung
ausgelöster Halt bestätigte die vorgesehene Anwendung des
\emph{fail-closed}-Prinzips. Technische Implementierungsfehler wurden
korrigiert und der betroffene Verarbeitungsschritt anschließend gegen
dieselben Akzeptanzkriterien erneut geprüft.

\paragraph{Ergebnisse der Verarbeitung mehrerer Quellen}

Der erste zusammenhängende Verarbeitungslauf mit mehreren Quellen machte
eine Identitätskollision zwischen Zuständen unterschiedlicher
Verarbeitungsläufe sichtbar. Technische Identifikatoren, die innerhalb eines
einzelnen Laufs eindeutig waren, konnten an einer projektweiten
Verarbeitungsgrenze nicht mehr eindeutig ihrem Quellen- und
Verarbeitungskontext zugeordnet werden. Die Weiterverarbeitung wurde an
dieser Stelle entsprechend dem \emph{fail-closed}-Prinzip angehalten.

Zur Korrektur wurde die projektweite Referenzierung um den jeweiligen
Verarbeitungskontext ergänzt. Dadurch konnten auch lokal identische
Kennungen eindeutig über ihren Verarbeitungslauf unterschieden und mit den
zugehörigen Quellen-, Prüf- und Informationszuständen verknüpft werden.
Anschließend wurde der Verarbeitungspfad erneut geprüft.

Im Wiederholungslauf konnten die Engineering-Quellen zunächst getrennt
verarbeitet und ihre autorisierten Engineering-Informationen mit der
jeweiligen Quellen- und Verarbeitungsprovenienz erhalten werden. Über die
projektbezogene Quellenzulassung wurde anschließend geprüft, welche Quellen
für die projektweite Verarbeitung berücksichtigt wurden. Die zugelassenen
Informationszustände konnten daraufhin gemeinsam für die Modellbildung
bereitgestellt werden.

Tabelle~\ref{tab:validation_multisource_results} fasst die im
Wiederholungslauf beobachteten Ergebnisse zusammen.

\begin{table}[htbp]
    \centering
    \small
    \caption[Ergebnisse der Verarbeitung mehrerer Quellen]{
        Ergebnisse der Verifikation der Verarbeitung mehrerer
        Engineering-Quellen nach Korrektur der Identitätskollision.}
    \label{tab:validation_multisource_results}
    \begin{tabular}{p{4.4cm}p{8.6cm}}
        \hline
        \textbf{Prüfgegenstand} &
        \textbf{Beobachtetes Ergebnis} \\
        \hline

        Quellenidentität &
        Getrennt registrierte Engineering-Quellen blieben während der
        quellenbezogenen Verarbeitung eindeutig unterscheidbar. \\

        Provenienz &
        Zustände aus unterschiedlichen Verarbeitungsläufen konnten nach der
        Korrektur eindeutig ihrem jeweiligen Quellen- und
        Verarbeitungskontext zugeordnet werden. \\

        Fachliche Autorisierung &
        Die quellenbezogen autorisierten Engineering-Informationen blieben
        als eigenständige referenzierbare Zustände erhalten. \\

        Projektbezogene Quellenzulassung &
        Die projektweite Verarbeitung berücksichtigte die gespeicherten
        Project-Fit-Zustände der beteiligten Quellen. \\

        Projektweite Übergabe &
        Die Übergabe erhielt die Bindungen an Quelle, Verarbeitungslauf,
        fachlich autorisierte Information und Quellenzulassung. \\

        Gemeinsame Modellbildung &
        Informationen mehrerer zugelassener Quellen konnten gemeinsam in
        die projektweite Modellbildung eingehen, während ihre jeweilige
        Herkunft und fachliche Autorisierung erhalten blieben. \\

        Gesamtergebnis &
        Der korrigierte Verarbeitungspfad mit mehreren Quellen konnte bis
        zur gemeinsamen projektweiten Modellbildung vollständig durchlaufen
        werden. \\
        \hline
    \end{tabular}
\end{table}

Der Wiederholungslauf bestätigte damit die eindeutige Zuordnung der
quellenbezogenen Informations- und Autorisierungszustände auch bei ihrer
projektweiten Verarbeitung. Die jeweilige Quellen- und
Verarbeitungsprovenienz blieb erhalten und die zugelassenen Informationen
konnten gemeinsam für die nachfolgende Modellbildung bereitgestellt werden.

Der weitergehende semantische Abgleich zwischen Informationen verschiedener
Quellen und eine daraus gegebenenfalls entstehende neue Engineering-Aussage
wurden in diesem Wiederholungslauf nicht vollständig durchlaufen. Die
Evaluation dieses Verarbeitungspfads konzentrierte sich auf die
quellenbezogene Verarbeitung und Autorisierung, die projektbezogene
Quellenzulassung sowie die gemeinsame Bereitstellung der zugelassenen
Informationen für die Modellbildung.

\subsection{Validierung des SysML-v2-Artefakts}
\label{sec:validation_sysml_human}

Die Validierung des erzeugten SysML-v2-Artefakts untersucht dessen formale
und strukturelle Eigenschaften sowie seine Verarbeitbarkeit in der
vorgesehenen SysML-v2-Zielumgebung. Prüfgegenstand ist dabei der konkret
erzeugte Artefaktzustand.

Zunächst führt die Implementierung interne regelbasierte Prüfungen durch.
Dabei werden insbesondere die Zuordnung des erzeugten Artefakts zum
zugrunde liegenden internen Modell- und Generierungszustand, die
Vollständigkeit der erforderlichen Metadaten sowie die für die weitere
Verarbeitung festgelegten Invarianten geprüft.

Anschließend wird das erzeugte Artefakt mit der
\emph{SYSIDE Modeler CLI} in einer externen SysML-v2-Toolumgebung
verarbeitet. Als Akzeptanzkriterien gelten die erfolgreiche Verarbeitung
des Artefakts und ein Validierungszustand ohne blockierende Diagnosen.
Damit wird geprüft, ob die erzeugte SysML-v2-Repräsentation auch außerhalb
der eigenen Generierungslogik formal verarbeitet werden kann.

Ergänzend wurde geprüft, ob der Prototyp die für die anschließende
menschliche Modellprüfung und Freigabe erforderlichen Zustände eindeutig
bereitstellt und miteinander verknüpft. Dazu gehören insbesondere der
erzeugte Artefaktzustand, das zugehörige Validierungsergebnis sowie die für
den Modellinhalt maßgeblichen fachlichen und modellbezogenen
Autorisierungsentscheidungen. Eine erfasste Modellprüfung und
Freigabeentscheidung kann dadurch eindeutig an den tatsächlich geprüften
Artefaktzustand gebunden werden.

\paragraph{Ergebnisse}

Das mit dem konsolidierten Prototyp erzeugte SysML-v2-Artefakt konnte mit
der \emph{SYSIDE Modeler CLI}~0.10.3 erfolgreich verarbeitet werden. Der
gespeicherte Validierungszustand weist den Status \texttt{valid}, einen
bestandenen \texttt{Publication gate}, keine Findings und keine
Diagnosemeldungen der externen Toolumgebung aus.

Auch die für die anschließende menschliche Modellprüfung und Freigabe
erforderlichen Zustände lagen eindeutig verknüpft vor. Das erzeugte
SysML-v2-Artefakt, sein Validierungsergebnis sowie die zugrunde liegenden
Informations-, Autorisierungs- und Modellzustände konnten demselben
Verarbeitungspfad zugeordnet werden. Die im Demonstrationslauf erfasste
Freigabeentscheidung wurde entsprechend an den konkret geprüften
Artefaktzustand gebunden.

Die vollständige Rekonstruktion dieser Verknüpfungen anhand des
abschließenden Demonstrationslaufs wird in
Kapitel~\ref{sec:results} dargestellt.

Die im Demonstrationslauf gespeicherten Informations-, Entscheidungs- und
Modellzustände bilden die Grundlage für die in
Kapitel~\ref{sec:results} dargestellte Rekonstruktion der durchgängigen
Nachverfolgbarkeit.